\section*{Chapter \ref{ch:dv:trading-volatility} --- Trading Volatility}

\begin{solution}{exo:dv:trading-volatility:1}
$767.5\times(0.02^2-0.18^2/252)=767.5\times0.000271=0.208$ for the 2\% day; $-767.5\times0.18^2/252=-0.099$ for the quiet day. The breakeven
daily move is $0.18/\sqrt{252}=1.13\%$.
\end{solution}

\begin{solution}{exo:dv:trading-volatility:2}
The realised variance came when the straddle's \omterm{def:dv:greeks-and-the-hedging-pnl:cashgamma}{cash gamma} was small. The quiet weeks (7.9\% realised) fell at the strike, where the \omterm{def:dv:greeks-and-the-hedging-pnl:cashgamma}{cash gamma} was
about 770 and the straddle paid \omterm{def:dv:greeks-and-the-hedging-pnl:greeks}{theta} for nothing: $-0.80$. The wild week (45\%) came with the share at 107.8, where the \omterm{def:dv:greeks-and-the-hedging-pnl:cashgamma}{cash gamma} was 39, a
twentieth of the starting level: $+0.02$.
\end{solution}

\begin{solution}{exo:dv:trading-volatility:3}
Carry $-1.88$: time decay at an unchanged 21.79 is $-1.59$ and roll-down $-0.28$ (the volatility falls 0.87 point to the two-month 20.92). On an
inverted curve the shorter expiry has the higher volatility, so a long position rolls up: positive roll-down, which offsets part of the decay.
\end{solution}

\begin{solution}{exo:dv:trading-volatility:4}
Short one-month straddles decay faster per unit of \omterm{def:dv:greeks-and-the-hedging-pnl:greeks}{vega} than the long three-month straddle, and the three-month leg rolls down less per day. With
the \omterm{def:dv:greeks-and-the-hedging-pnl:greeks}{vega} matched (1.73 one-month straddles per three-month straddle) the net \omterm{def:dv:greeks-and-the-hedging-pnl:greeks}{theta} is positive: 0.13 a day. The price is a short gamma of $-0.17$, so
a large move in the next days loses money, and a flattening of the curve does too.
\end{solution}

\begin{solution}{exo:dv:trading-volatility:5}
$\sigma(K)=a+s\ln(K/S)/\sqrt\tau$. At fixed $K$ and $a$, $\sigma_1-\sigma_0=s\bigl(\ln(K/S_1)-\ln(K/S_0)\bigr)/\sqrt\tau=-s\ln(S_1/S_0)/\sqrt\tau$, the same for
every strike. With $s=-0.10$, $S_1/S_0=0.95$ and $\tau=0.25$: $0.10\times\ln0.95/0.5=-0.0103$, a fall of about one point.
\end{solution}

\begin{solution}{exo:dv:trading-volatility:6}
Gamma and \omterm{def:dv:greeks-and-the-hedging-pnl:greeks}{theta} depend on the path of the spot and on the marks at the start of each day, which are close under the two rules when the spot is near
its reference. \omterm{def:dv:greeks-and-the-hedging-pnl:greeks}{Vega} and \omterm{def:dv:greeks-and-the-hedging-pnl:greeks}{vanna} depend on $\delta\sigma_i$, the day's change in each option's mark. Under \omterm{def:dv:implied-volatility-and-its-surface:sticky}{sticky strike} it comes from the level only,
under \omterm{def:dv:implied-volatility-and-its-surface:sticky}{sticky delta} also from the spot's move, so the two rules assign the same move to different columns.
\end{solution}

\begin{solution}{exo:dv:trading-volatility:7}
Without a premium the average is $-0.008\%$ of the spot a month, with a standard deviation of 1.33, so a standard error of $1.33/\sqrt{120}=0.12$: indistinguishable
from zero. The Sharpe ratio is $-0.02$. The implied volatility sold (21.5 on average) exceeds the average realised (20.6) because it is the square root of the
expected variance, and by Jensen's inequality that exceeds the expected square root. In variance the two match: root-mean-squares of 23.0 and 23.2. The two
points of premium are the whole of the $+0.44\%$.
\end{solution}

\begin{solution}{exo:dv:trading-volatility:8}
The Sharpe ratio and the hit rate describe the body of the distribution, not the tail. In the chapter's programme the same figures (1.15, 74\%) coexist with a
worst month of $-5.04\%$, 4.1 standard deviations below the mean, and a skewness of $-1.35$. Ten years may contain no crash of the size that sets the true risk. Measure the
risk by stress (a 5 February 2018) and by tail measures, not by volatility.
\end{solution}

\begin{solution}{pb:dv:trading-volatility:1}
\textbf{1.} 4.15; that the gamma-weighted realised variance over the month would exceed 18\% volatility.
\textbf{2.} $\pm1.57\%$ a day; $+3.28$.
\textbf{3.} About 770 ($\frac12\Gamma S^2$); a daily move of $0.18/\sqrt{252}=1.13\%$ breaks even.
\textbf{4.} 3.28 against 3.05: the higher-order terms of daily moves of 1.6\%, which the second-order sum leaves out.
\textbf{5.} The same on all three paths: the notional times $25^2-18^2$ in variance points, since a \omterm{def:dv:variance-swaps-and-volatility-derivatives:vs}{variance swap}'s \omterm{def:dv:greeks-and-the-hedging-pnl:cashgamma}{cash gamma} is constant.
\textbf{6.} Fifteen days of $+0.5\%$ and six of $\pm2.84\%$: $15\times0.005^2+6\times0.0284^2=0.00521=0.25^2\times21/252$.
\textbf{7.} $-0.80$ in the quiet weeks, $+0.02$ in the wild week: $-0.78$.
\textbf{8.} At about 107.8, 7.8\% above the strike: a \omterm{def:dv:greeks-and-the-hedging-pnl:cashgamma}{cash gamma} of 39 against 770 at inception.
\textbf{9.} 2.87. Its quiet weeks, at the strike, pay \omterm{def:dv:greeks-and-the-hedging-pnl:greeks}{theta} at full gamma with 7.9\% realised; the even path earns on every day.
\textbf{10.} The spread is the expected P\&L only on average over paths. On one path the realised variance's timing relative to the gamma decides; a \omterm{def:dv:variance-swaps-and-volatility-derivatives:vs}{variance swap}
isolates the spread.
\textbf{11.} Re-striking at the money (selling the old straddle, buying a new one) would have kept the gamma near its starting level for the wild week, at the cost of
the bid--ask spread and a new premium.
\textbf{12.} $-1.88$ over the month, of which $-0.28$ is roll-down.
\textbf{13.} 0.13 a day, for a short gamma of $-0.17$: a large move or a flattening curve.
\textbf{14.} $-0.023$ marked \omterm{def:dv:implied-volatility-and-its-surface:sticky}{sticky strike}, $+0.081$ \omterm{def:dv:implied-volatility-and-its-surface:sticky}{sticky delta}.
\textbf{15.} The Greek explain misses 0.08 either way: third-order terms of a 5\% move are as large as the day's P\&L.
\textbf{16.} Not about the level: the share did realise 25. The trade was a bet on a path-weighted quantity, and the path went against it.
\textbf{17.} Gamma earned on the quiet weeks below \omterm{def:dv:greeks-and-the-hedging-pnl:greeks}{theta}, almost no gamma in the wild week, no \omterm{def:dv:greeks-and-the-hedging-pnl:greeks}{vega} (the implied volatility did not change), and a small unexplained
from the large moves.
\textbf{18.} Many small gains and a few large losses (skewness $-1.35$, worst month $-5.04\%$): the premium pays for bearing the crashes.
\textbf{19.} With the same 25\% realised, the long straddle loses 0.78 when the moves come on the low-gamma days and makes 3.28 when they are spread evenly.
\textbf{20.} The realised variance weighted by the option's \omterm{def:dv:greeks-and-the-hedging-pnl:cashgamma}{cash gamma}, minus the implied variance on the same weights.
\end{solution}

\begin{solution}{iq:dv:trading-volatility:1}
From the difference between the realised squared moves and the implied variance, weighted each day by the \omterm{def:dv:greeks-and-the-hedging-pnl:cashgamma}{cash gamma} $\frac12\Gamma S^2$: gamma gains when the
moves are large, \omterm{def:dv:greeks-and-the-hedging-pnl:greeks}{theta} losses when they are small; plus \omterm{def:dv:greeks-and-the-hedging-pnl:greeks}{vega} when the implied volatility moves, and the cross terms.
\iqlookfor{the cash-gamma-weighted spread.}
\end{solution}

\begin{solution}{iq:dv:trading-volatility:2}
The realised variance came where the gamma was small: the share drifted away from the strike quietly and moved sharply far from it, or the big moves came when
the option had much time left and little gamma per unit of move. Also possible: the implied volatility fell (\omterm{def:dv:greeks-and-the-hedging-pnl:greeks}{vega} loss), or hedging at discrete times missed
intraday moves.
\iqlookfor{path dependence through the gamma, and \omterm{def:dv:greeks-and-the-hedging-pnl:greeks}{vega}.}
\end{solution}

\begin{solution}{iq:dv:trading-volatility:3}
Long the longer expiry, short the shorter, vega-neutral: the short leg's faster decay gives positive \omterm{def:dv:greeks-and-the-hedging-pnl:greeks}{theta}, and the long leg rolls down the curve, which costs a
little. Net carry is usually positive; the risk is short gamma and a flattening or inversion of the curve.
\iqlookfor{\omterm{def:dv:greeks-and-the-hedging-pnl:greeks}{theta}, roll-down, and the short gamma that pays for them.}
\end{solution}

\begin{solution}{iq:dv:trading-volatility:4}
Per option, at its own mark: delta, gamma, \omterm{def:dv:greeks-and-the-hedging-pnl:greeks}{theta}, \omterm{def:dv:greeks-and-the-hedging-pnl:greeks}{vega}, \omterm{def:dv:greeks-and-the-hedging-pnl:greeks}{vanna}, \omterm{def:dv:greeks-and-the-hedging-pnl:greeks}{volga} from the start-of-day \omterm{def:dv:greeks-and-the-hedging-pnl:greeks}{Greeks} and the day's moves in spot and in each mark; full revaluation
alongside. The unexplained part is monitored: small and random is fine, while large or persistent means missing risk factors, large moves, a marking change or
a booking error, each to be investigated.
\iqlookfor{per-option explain, full revaluation, unexplained as a control.}
\end{solution}

\begin{solution}{iq:dv:trading-volatility:5}
Neither is right in general: they are two assumptions about how the surface moves with the spot. The total P\&L of the day follows the marks, and the split
follows the rule. Choose the rule the market has followed for this underlying (chapter 7), and report the \omterm{def:dv:greeks-and-the-hedging-pnl:greeks}{vanna} exposure it creates as a risk.
\iqlookfor{marking rule as an assumption, and its effect on the split.}
\end{solution}

\begin{solution}{iq:dv:trading-volatility:6}
No: it is payment for bearing crash risk and for providing insurance. Returns have a long left tail. Size by stress loss (a 5 February 2018 or worse), not by
volatility; cap the short gamma and \omterm{def:dv:greeks-and-the-hedging-pnl:greeks}{vega}; prefer defined-risk structures; keep capital for the tail.
\iqlookfor{compensation for tail risk, sizing by stress.}
\end{solution}
